% Part:sets-functions-relations
% Chapter: size-of-sets
% Section: enumerations-alt

\documentclass[../../../include/open-logic-section]{subfiles}

\begin{document}

\olfileid[fr]{sfr}{siz}{enm-alt}

\olsection{Énumérations bijectives et ensembles \usetoken{p}{enumerable}}

\begin{editorial}
  Cette section définit les énumérations comme des bijections avec
  $\Nat$ ou avec un segment initial de~$\Nat$, selon l'usage de
  la théorie des ensembles. Cette définition diffère donc légèrement
  de celle de la \olref[enm]{sec}, dont tous les exemples sont repris
  ici. La présentation est aussi un peu plus concise.
\end{editorial}

On peut décrire un ensemble fini en énumérant simplement ses éléments,
comme lorsqu'on écrit :
\[
  A = \{a_1, a_2, \ldots, a_n\}.
\]
Si les éléments $a_1$, \dots, $a_n$ sont tous distincts, cette
écriture donne une !!{bijection} entre $A$ et les $n$ premiers
entiers naturels $0$, \dots, $n-1$. Réciproquement, tout ensemble
fini ne possède qu'un nombre fini d'éléments et peut donc être mis
en une telle correspondance. Autrement dit, si $A$ est fini, il
existe une !!{bijection} entre $A$ et $\{0, \dots, n-1\}$, où
$n$ est le nombre d'éléments de~$A$.\footnote{Précision éditoriale :
pour $n=0$, le segment formé des $n$ premiers entiers est vide.
La définition d'énumération ci-dessous utilise quant à elle des
segments non vides ; elle traite l'ensemble vide séparément dans
la définition de la dénombrabilité.}

Si l'on admet certains ensembles infinis, on peut aussi admettre
des énumérations infinies. Précisons cette idée : un ensemble infini
est énuméré par une !!{bijection} entre cet ensemble et tout~$\Nat$.

\begin{defn}[Énumération au sens ensembliste]
Une \emph{énumération} d'un ensemble $A$ est une !!{bijection} dont
l'image est $A$ et dont le domaine est soit un segment initial
$\{0, 1, \ldots, n\}$ des entiers naturels, soit l'ensemble
$\Nat$ tout entier.
\end{defn}

\begin{explain}
L'emploi du mot \emph{énumération} repose sur une idée intuitive :
énumérer un ensemble $A$, c'est disposer d'une !!{bijection} $f$
qui permet d'en compter les éléments un à un. L'élément d'indice
$0$ est $f(0)$, celui d'indice $1$ est $f(1)$, \ldots, celui
d'indice $n$ est $f(n)$, \ldots.\footnote{Nous commençons bien
à compter à partir de $0$. On pourrait aussi commencer à $1$ :
cela ne changerait rien d'essentiel. Il suffirait de remplacer
$\Nat$ par~$\PosInt$.} La définition suivante précise encore
cette idée.
\end{explain}

\begin{defn}
  \ollabel{defn:enumerable}
  Un ensemble~$A$ est !!{enumerable} si et seulement si
  $A = \emptyset$ ou s'il existe une énumération de~$A$.
  On dit que $A$ est !!{nonenumerable} si et seulement s'il
  n'est pas !!{enumerable}.
\end{defn}

\begin{explain}
Ainsi, un ensemble est !!{enumerable} si et seulement s'il est
vide ou si une énumération permet d'en compter les éléments un à un.
\end{explain}

\begin{ex}
L'identité $\Id{\Nat} \colon \Nat \to \Nat$, définie par
$\Id{\Nat}(n) = n$, énumère les entiers naturels. La fonction
$g(n) = n + 1$, c'est-à-dire la fonction successeur, énumère les
entiers naturels \emph{strictement positifs}
$\Nat^+ = \Nat \setminus \{0\}$.
\end{ex}

\begin{prob}
Montrer qu'un ensemble $A$ est !!{enumerable} si et seulement si
$A = \emptyset$ ou s'il existe une !!{surjection}
$f\colon \Nat \to A$. Montrer aussi que $A$ est !!{enumerable}
si et seulement s'il existe une !!{injection} $g\colon A \to \Nat$.
\end{prob}

\begin{ex}
Les fonctions $f\colon \Nat \to \Nat$ et $g \colon \Nat \to \Nat$
définies par
\begin{align*}
  f(n) & = 2n \text{ et}\\
  g(n) & = 2n+1
\end{align*}
énumèrent respectivement les entiers naturels pairs et les entiers
naturels impairs, lorsque l'on prend leur image respective comme
ensemble d'arrivée.\footnote{Précision éditoriale : ces applications
deviennent des bijections par coréstriction à leur image.
L'original les appelle directement énumérations tout en les écrivant
à valeurs dans $\Nat$ ; cette distinction est nécessaire avec
la définition bijective adoptée ici.}
Mais aucune des deux fonctions à valeurs dans $\Nat$ n'est
!!{surjective} ; aucune n'est donc une énumération de~$\Nat$.
\end{ex}

\begin{prob}
Définir une énumération des nombres carrés $1$, $4$, $9$, $16$, \dots
\end{prob}

\begin{ex}
Notons $\lceil x \rceil$ le \emph{plafond} de $x$, c'est-à-dire
le plus petit entier supérieur ou égal à~$x$. La fonction
$f \colon \Nat \to \Int$ définie par
\[
  f(n) = (-1)^{n} \left\lceil\tfrac{n}{2}\right\rceil
\]
énumère l'ensemble des entiers relatifs~$\Int$ de la manière suivante :
\[
\begin{array}{c c c c c c c c}
f(0) & f(1) & f(2) & f(3) & f(4) & f(5) & f(6) & \dots \\ \\
\big\lceil \tfrac{0}{2} \big\rceil & -\big\lceil \tfrac{1}{2}\big\rceil &  \big\lceil \tfrac{2}{2} \big\rceil & -\big\lceil \tfrac{3}{2} \big\rceil & \big\lceil \tfrac{4}{2} \big\rceil  & -\big\lceil \tfrac{5}{2}\big\rceil & \big\lceil \tfrac{6}{2} \big\rceil & \dots \\ \\
0 & -1 & 1 & -2 & 2 & -3 & 3& \dots
\end{array}
\]
La fonction $f$ produit les valeurs de $\Int$ en « sautant »
d'un côté à l'autre de zéro, entre entiers positifs et négatifs.
On peut aussi la définir par cas :
\[
f(n) = \begin{cases}
  \frac{n}{2} & \text{si $n$ est pair}\\
  -\frac{n+1}{2} & \text{si $n$ est impair}
  \end{cases}
\]
\end{ex}

\begin{prob}
Montrer que, si $A$ et $B$ sont !!{enumerable}s, alors
$A \cup B$ l'est aussi.
\end{prob}

\begin{prob}
Montrer par récurrence sur $n$ que, si $A_1$, $A_2$, \dots,
$A_n$ sont tous !!{enumerable}s, alors $A_1 \cup \dots \cup A_n$
l'est aussi.
\end{prob}

\end{document}
